\documentclass[12pt,a4paper]{article}
\usepackage[utf8]{inputenc}
\usepackage[T1]{fontenc}
\usepackage[italian]{babel}
\usepackage{geometry}
\usepackage{amsmath}
\usepackage{hyperref}
\usepackage{enumitem}
\usepackage{booktabs}
\usepackage{parskip}

\geometry{margin=2.5cm}

\hypersetup{colorlinks=true, linkcolor=blue, urlcolor=blue}

\title{\textbf{Progetto Esteso -- TLN Parte III}\\
\large RAG Ibrido con Valutazione Quantitativa\\
\normalsize Riassunto per lo studio}
\author{Tecnologie del Linguaggio Naturale -- Prof.\ Di Caro -- A.A.\ 2025/2026}
\date{}

\begin{document}

\maketitle
\tableofcontents
\newpage

% ─────────────────────────────────────────────
\section{Cos'è RAG --- la base del progetto}

RAG (\textbf{Retrieval-Augmented Generation}) è un sistema che risponde a domande
combinando due componenti:
\begin{enumerate}
  \item Un \textbf{motore di ricerca} che recupera i documenti più rilevanti per la domanda;
  \item Un \textbf{LLM} (Large Language Model) che legge quei documenti e genera la risposta.
\end{enumerate}

Il vantaggio è che il modello non deve ``ricordare'' tutto: accede a conoscenza
esterna al momento del bisogno. Il flusso è:
\[
  \text{Query utente}
  \xrightarrow{\text{retrieval}}
  \text{Documenti rilevanti}
  \xrightarrow{\text{generazione}}
  \text{Risposta}
\]

\subsection{Cosa abbiamo costruito nel LAB-5}

Nel quinto laboratorio abbiamo implementato un sistema RAG completo sul dataset
\texttt{MaartenGr/arxiv\_nlp} di Hugging Face, contenente circa 45.000 abstract
di paper scientifici NLP. Il sistema:
\begin{itemize}
  \item genera \textbf{embeddings semantici} degli abstract con \texttt{all-MiniLM-L6-v2};
  \item indicizza i vettori con \textbf{FAISS} (Fast Approximate Nearest Neighbor);
  \item recupera i $k$ documenti più simili alla query;
  \item genera la risposta con un LLM locale (es.\ Qwen2.5-0.5B).
\end{itemize}

% ─────────────────────────────────────────────
\section{Il problema del RAG base: perché non basta}

Il retrieval basato solo su embeddings ragiona per \textbf{concetti}, non per
\textbf{parole esatte}. Questo crea problemi in due casi tipici:
\begin{itemize}
  \item Query con \textbf{termini tecnici specifici} (nomi di algoritmi, acronimi);
  \item Query \textbf{molto brevi} dove il contesto semantico è insufficiente.
\end{itemize}

\textbf{Esempio.} Se cerchi ``BERT pre-training'', il sistema semantico potrebbe
restituire documenti che parlano di transformer in generale, anche se non
contengono letteralmente quelle parole.

% ─────────────────────────────────────────────
\section{Il progetto esteso: Hybrid Search}

\subsection{BM25 --- il retrieval lessicale}

BM25 (\textit{Best Match 25}) è un algoritmo di ranking classico che misura la
rilevanza di un documento rispetto a una query basandosi sulla
\textbf{frequenza delle parole} nel documento e nella collezione.
A differenza degli embeddings, BM25 lavora sulle \textbf{parole esatte}.

La formula di scoring è:
\[
  \text{score}(D,Q) = \sum_{i=1}^{n} \mathrm{IDF}(q_i)
  \cdot \frac{f(q_i,D)\cdot(k_1+1)}
             {f(q_i,D)+k_1\!\left(1-b+b\cdot\dfrac{|D|}{\mathrm{avgdl}}\right)}
\]
dove $f(q_i,D)$ è la frequenza del termine nel documento, $|D|$ è la lunghezza
del documento, \textit{avgdl} è la lunghezza media, e $k_1=1.5$, $b=0.75$ sono
parametri tipici.

\subsection{Confronto tra i due approcci}

\begin{center}
\begin{tabular}{lll}
\toprule
& \textbf{Embeddings semantici} & \textbf{BM25 lessicale} \\
\midrule
Punto di forza   & Capisce sinonimi e concetti    & Preciso su termini esatti \\
Punto debole     & Impreciso su termini specifici & Non capisce il significato \\
Buono per        & Query concettuali e lunghe     & Query con termini tecnici \\
\bottomrule
\end{tabular}
\end{center}

\subsection{Hybrid Search: la combinazione}

Combiniamo i due metodi con \textbf{Reciprocal Rank Fusion (RRF)}:
\[
  \mathrm{RRF}(d) = \sum_{r \in R} \frac{1}{k + r(d)}
\]
dove $r(d)$ è il rank del documento $d$ nella lista $r$ e $k=60$ è una costante
tipica. In pratica:
\begin{enumerate}
  \item BM25 produce una lista ordinata di documenti;
  \item Gli embeddings producono un'altra lista ordinata;
  \item RRF fonde le due liste dando peso ai documenti che appaiono in cima a entrambe.
\end{enumerate}

% ─────────────────────────────────────────────
\section{La valutazione quantitativa}

Non ci limitiamo a dire ``funziona bene'': \textbf{misuriamo} quanto funziona bene.

\textbf{Precision@k} -- dei primi $k$ documenti restituiti, quanti sono rilevanti?
\[
  P@k = \frac{|\text{documenti rilevanti tra i primi } k|}{k}
\]

\textbf{Recall@k} -- di tutti i documenti rilevanti esistenti, quanti ne ho recuperati?
\[
  R@k = \frac{|\text{documenti rilevanti tra i primi } k|}{|\text{tutti i documenti rilevanti}|}
\]

\textbf{Faithfulness} -- la risposta generata è fedele ai documenti recuperati,
o il modello ``inventa'' informazioni non presenti nel contesto?

\subsection{Disegno sperimentale}

\begin{center}
\begin{tabular}{lll}
\toprule
\textbf{Configurazione} & \textbf{Retrieval} & \textbf{Atteso migliore su} \\
\midrule
RAG base    & Solo embeddings         & Query concettuali \\
RAG ibrido  & BM25 + embeddings (RRF) & Query con termini precisi \\
\bottomrule
\end{tabular}
\end{center}

% ─────────────────────────────────────────────
\section{Stack tecnologico}

\begin{itemize}
  \item \texttt{rank\_bm25} -- implementazione Python di BM25
  \item \texttt{sentence-transformers} -- generazione embeddings
  \item \texttt{faiss-cpu} -- vector store per nearest neighbor search
  \item \texttt{datasets} (HuggingFace) -- caricamento dataset arXiv NLP
  \item \texttt{transformers} + \texttt{torch} -- LLM locale per la generazione
  \item \texttt{RAGAS} (opzionale) -- framework per valutazione automatica RAG
\end{itemize}

% ─────────────────────────────────────────────
\section{Cosa dire al professore}

\textit{``Abbiamo completato tutti e cinque i laboratori. Come progetto esteso
vorremmo partire dal LAB-5 e estenderlo in due direzioni.}

\textit{La prima è implementare un hybrid search: invece di usare solo gli
embeddings semantici, combiniamo BM25 con gli embeddings tramite Reciprocal
Rank Fusion, e confrontiamo i risultati su query di tipo diverso.}

\textit{La seconda è aggiungere una valutazione quantitativa con Precision@k,
Recall@k e, se riusciamo, la faithfulness.}

\textit{L'obiettivo non è solo far funzionare il sistema, ma confrontare almeno
due configurazioni e analizzare i trade-off.''}

% ─────────────────────────────────────────────
\section{Connessione con il corso}

\begin{itemize}
  \item \textbf{Cap.\ 6 -- Pre-LLM NLP}: BM25 è citato esplicitamente come
        baseline robusta per il retrieval lessicale.
  \item \textbf{Cap.\ 7 -- RAG}: il progetto è un'estensione diretta del LAB-5,
        citata nella sezione 7.14.8 della dispensa.
  \item \textbf{Cap.\ 9 -- Valutazione}: le metriche Precision@k e faithfulness
        sono al centro della griglia di valutazione del prof.
\end{itemize}

\end{document}